\scanpage{27}
% lines 1–7
En car. 2, la catégorie des formes de $\mathfrak{sl}(2)$ est équivalente à celle des couples $(V,E)$ où $V$ est un Module vectoriel de rang 2 et $E$ une extension de \deleted{$\unsure{E}$} $V$ par $\det V$ (\uline{NB} pour $V$ fixé, la donnée d'une telle extension équivaut à celle d'une \deleted{\unsure{fibré}} extension de $\det V^{-1}$ par $\unsure{\check V^*|\check V}$ ou de $\mathscr O_S$ par $(\det V)\otimes\check V\simeq V$, i.e. à la cat. des fibrés affines sous $V\ldots$).

% lines 7–12
La catégorie des formes de $\mathfrak{gp}(1)$ équivaut à celle des \deleted{\unsure{extens}} couples $(V',E')$ d'un Module vectoriel \inserted{$V'$} de rang 2, et d'une extension \inserted{$E'$} de $\mathscr O_S$ par $V'$ ; (pour $V'$ fixé, la cat équivaut à celle des fibrés affines sous $V'$).

% lines 12–14
Disons que deux telles algèbres de Lie \inserted{$E,E'$} sont accouplées si on a un accouplement dualisant entre $E,E'$, dans

% lines 15–16
\diagram{\begin{tikzcd}[ampersand replacement=\&,column sep=small,row sep=small]
\& \mathscr O_S \arrow[d,phantom,"\rotatebox{90}{$\simeq$}" description] \& \& \& \\
\unsure{\bullet|0} \arrow[r,no head] \& \det V \arrow[r,no head] \& E \arrow[r] \& V \arrow[r] \& 0\\
\bullet \arrow[r,no head] \& \deleted{\unsure{\det}}\mathscr O_S \& E' \arrow[l] \& \deleted{\unsure{\check{\phantom V}}}V' \arrow[l] \arrow[d,phantom,"\rotatebox{90}{$\simeq$}" description] \& \bullet \arrow[l]\\
\& \& \& V \&
\end{tikzcd}}

% lines 15–20
lequel le sous-module $\det V$ de $E$ est l'orthogonal du sous module $V'$ de $E'$ -- qui est donc \deleted{iso à $\check V$ =} dual de $V$, et de = $\det V$ est dual de $\mathscr O_S$, donc can. iso. à $\mathscr O_S$. \deleted{Donc} $V$ est en fait muni d'une structure symplectique, \deleted{et} \inserted{par laquelle} $V'\simeq V$.

% lines 20–26
La catégorie des couples \inserted{accouplés} de formes de $\mathfrak{sl}(2),\mathfrak{gp}(1)$ est donc équivalente à celle des triples $(V,\varphi,E)$ d'un Module vectoriel $V$ de rang 2, avec structure symplectique $\varphi$ et ext. de $V$ par $\mathscr O_S$ (ce qui revient à ext. de $\mathscr O_S$ par $\check V\simeq V$, ou encore fibré affine sous $V$) : c'est donc équivalent

% lines foot
% archive sheet number: 26
